% Chapter introduction: why the topic arose, who first asked, what it answers.
% STATUS: rewritten 2026-11-05; every claim is cited in its section.

The earlier chapters studied derivations on rings, algebras, algebras of
operators and Lie type structures. This chapter asks how the product rule
behaves when the ambient system loses part of the ring structure. A
near-ring drops one distributive law. A semiring drops the additive
negatives. A $\Gamma$-ring replaces the single product by a family of
products. A lattice keeps only the order. A BCI, MV or subtraction algebra
keeps only an implication or a difference.

\hbfigure{22-other-systems/structure-loss}{What each system in this chapter
keeps of the ring axioms, and which axiom it gives up.}{fig:other:structure}

Kuratowski asked the question first, but in the language of topology. His
closure and derivation operators on parts of a space were maps on a
Boolean algebra. Goldberg recalls them in this form in 1975
\cite[p.~365]{Goldberg1975}. The derived set of a set $A$ is the set of
its limit points. The operator $A \mapsto A^{d}$ is then a derivation on
the Boolean algebra of all parts of the space. This is the oldest
derivation on a system with no ring addition.

Bell and Mason moved the question to near-rings in 1987. A near-ring has
the product rule but only one distributive law. The proofs from prime
rings do not carry over. Their note gave the first commutativity theorems
for near-rings with a derivation
\cite[pp.~31--33]{BellMason1987}. Section~\ref{sec:other-systems:near-rings}
opens with these.

Between rings and lattices sits another family. Nobusawa $\Gamma$-rings
replace the product $xy$ by a product $x\, \alpha\, y$ indexed by a set
$\Gamma$ of scalars. Dey, Paul and Rakhimov treat derivations and
generalized derivations in this setting
\cite[pp.~2--3]{DeyPaulRakhimov2012}.
Section~\ref{sec:other-systems:gamma-rings} follows their definitions.

Semirings drop the negatives, so the proofs that use subtraction break. Ahmed
and Dudek show how left Jordan derivations force commutativity on a class
of semirings with an additive-inverse structure
\cite[pp.~624--631]{AhmedDudek2021}.
Section~\ref{sec:other-systems:semirings} takes this up.

Lattices come next. Xin, Li and Lu defined a derivation on a lattice by
the rule $d(x \wedge y) = (dx \wedge y) \vee (x \wedge dy)$
\cite[Definition~3.1, p.~309]{XinLiLu2008}. {\c{C}}even and {\"O}zt{\"u}rk
then replaced the second argument by a second map $f$ and obtained the
$f$-derivations \cite[Definition~8, p.~702]{CevenOzturk2008}. A residuated
lattice carries a product and its residuum, so He, Xin and Zhan could
define a multiplicative derivation there
\cite[Definition~3.1, p.~100]{HeXinZhan2016}.
Section~\ref{sec:other-systems:lattices} tells this part.

The last two sections treat algebras of logic. Jun and Xin applied the
derivation idea to BCI-algebras
\cite[Definition~3.5, p.~170]{JunXin2004}. Ciungu defined two types of
implicative derivations on pseudo-BCK algebras
\cite[p.~178]{Ciungu2021}. Wang, She and Qian studied implicative
derivations on MV-algebras
\cite[Definition~3.1, p.~262]{WangSheQian2019}. the subtraction of Schein
algebras, studied by {\"O}zt{\"u}rk and {\c{C}}even
\cite[Definition~3.1, p.~511]{OzturkCeven2009}, give the last family.
Sections~\ref{sec:other-systems:bci-mv}
and~\ref{sec:other-systems:subtraction-algebras} close the chapter.
